\chapter{Anaphylaxis}

\section*{Definition}
Anaphylaxis is a rapid, potentially fatal systemic reaction. It often begins within minutes to a few hours of exposure to a food, medicine, insect venom, or other trigger. It may involve the skin, breathing, circulation, or gastrointestinal tract, but skin symptoms can be absent. Because it can progress to airway obstruction and shock within minutes, treatment is based on the clinical picture and must not wait for laboratory testing.

\section*{When to See a Doctor}
Use your epinephrine auto-injector and call emergency services immediately if any of these develop after exposure to a possible allergen:
\begin{itemize}
  \item Trouble breathing, wheezing, throat tightness, hoarse voice, or feeling that your throat is closing
  \item Swelling of the lips, tongue, or face
  \item Dizziness, fainting, rapid weak pulse, or a sense of doom
  \item Repeated vomiting, severe cramping, or diarrhea together with symptoms in another body system
  \item Widespread hives or flushing plus any breathing, circulatory, or severe gastrointestinal symptom
  \item Sudden faintness, collapse, or low blood pressure after exposure to a known or likely allergen, even without hives
\end{itemize}
Do not wait for symptoms to worsen, and do not delay calling for help while you drive yourself -- have someone drive or call an ambulance.

\section*{How It Develops}
When an immune response is triggered, mast cells and basophils can release histamine and other mediators. These chemicals widen blood vessels and increase leakage into tissues, which can cause hives, swelling, and dangerously low blood pressure. They can narrow the airways and increase mucus, causing wheezing and breathing difficulty, and affect the gastrointestinal tract, causing vomiting or cramping. The reaction can involve several systems at once, but anaphylaxis can also occur without visible skin changes.

\section*{Causes}
\begin{itemize}
  \item Foods: most commonly peanuts, tree nuts, milk, eggs, fish, shellfish, wheat, and soy
  \item Medications: antibiotics such as penicillin, aspirin and other NSAIDs, and chemotherapy or contrast agents
  \item Insect venom: bee, wasp, hornet, and fire ant stings
  \item Latex and some vaccines
  \item Exercise-induced anaphylaxis, sometimes triggered when exercise follows a specific food
  \item Idiopathic anaphylaxis, where no trigger can be identified despite a careful workup
\end{itemize}

\section*{Related Symptoms}
\begin{itemize}
  \item Hives (raised, itchy welts) and flushed, warm skin
  \item Swelling of the lips, tongue, eyelids, or throat
  \item Wheezing, cough, chest tightness, and shortness of breath
  \item Nausea, vomiting, cramping, and diarrhea
  \item Dizziness, lightheadedness, and fainting from falling blood pressure
  \item Anxiety, a pounding heart, and a feeling of impending doom
\end{itemize}

\section*{Medical Evaluation}
\begin{itemize}
  \item The diagnosis is clinical, based on the timing and pattern of symptoms; no single blood test can confirm or exclude anaphylaxis during the emergency.
  \item During the reaction, clinicians repeatedly assess the airway, breathing, circulation, mental status, blood pressure, oxygen saturation, and response to treatment. Serum tryptase may be collected in selected cases, but it must never delay epinephrine.
  \item After recovery, an allergy-focused history and specialist testing (skin-prick testing or specific IgE blood tests) may help identify a trigger. A positive test shows sensitization and is not automatically proof of clinical allergy.
  \item For unclear triggers, supervised food or drug challenges may be arranged by an allergy specialist; they must not be attempted at home.
  \item The history is key: what was eaten, given, injected, or used; the timing; symptoms; treatment; and any previous reactions.
\end{itemize}

\section*{Treatment Options}
\begin{itemize}
  \item Epinephrine injected into the muscle (usually the outer thigh) is the first-line treatment and should be given as early as possible.
  \item Oxygen when hypoxemia or severe breathing difficulty is present, with close monitoring of breathing and circulation in the emergency department.
  \item Intravenous fluids to support blood pressure and reverse shock.
  \item A non-sedating antihistamine may help persistent hives or itching after epinephrine. Antihistamines and corticosteroids do not treat airway obstruction or shock and should never delay epinephrine; corticosteroids are not routinely used to prevent biphasic reactions.
  \item Bronchodilators such as albuterol for wheezing that persists after epinephrine.
  \item Observation for a period based on the severity of the reaction, treatment required, risk of recurrence, and local protocol, because symptoms can recur after initial improvement.
  \item Long-term prevention with immunotherapy (allergy shots) for insect-venom allergy, and strict avoidance of confirmed food and drug triggers.
\end{itemize}


\section*{Self-Care and Home Management}
\begin{itemize}
  \item Inject epinephrine in the outer thigh immediately when you suspect anaphylaxis -- do not wait to see how bad it gets.
  \item Lie flat with legs raised if this is comfortable and does not worsen breathing or vomiting. If breathing is difficult, sit with legs extended; if unconscious but breathing, use the recovery position. Do not stand or walk suddenly.
  \item Call emergency services and tell them anaphylaxis and the time of the epinephrine injection.
  \item Carry the number of epinephrine devices prescribed for you; many people at risk are advised to have two available because a second dose may be needed or the first device may fail. Follow your personal action plan.
  \item Antihistamines do not stop anaphylaxis -- they are not a substitute for epinephrine.
  \item After recovery, ask for a referral to an allergist and a written anaphylaxis action plan.
\end{itemize}

\section*{References}
\begin{thebibliography}{99}
\bibitem{anaphylaxis:wao2020} Cardona V, Ansotegui IJ, Ebisawa M, et al. World Allergy Organization Anaphylaxis Guidance 2020. \textit{World Allergy Organ J}. 2020;13:100472.
\bibitem{anaphylaxis:second-symposium} Sampson HA, Mu\u00f1oz-Furlong A, Campbell RL, et al. Second symposium on the definition and management of anaphylaxis: summary report. \textit{J Allergy Clin Immunol}. 2006;117:391--397.
\bibitem{anaphylaxis:wao2011} Simons FER, Ardusso LRF, Bil\u00f2 MB, et al. World Allergy Organization guidelines for the assessment and management of anaphylaxis. \textit{J Allergy Clin Immunol}. 2011;127:587--593.
\bibitem{anaphylaxis:practice-parameter} Lieberman P, Nicklas RA, Randolph C, et al. Anaphylaxis -- a practice parameter update 2015. \textit{Ann Allergy Asthma Immunol}. 2015;115:341--384.
\bibitem{anaphylaxis:practice-parameter-2023} Shaker MS, Wallace DV, Golden DBK, et al. Anaphylaxis: a 2023 practice parameter update. \textit{Ann Allergy Asthma Immunol}. 2023;130(1):1--63.
\end{thebibliography}
